\label{chap:prelim}
This chapter fixes the notation and the experimental ground on which the rest of the thesis stands: the navigation
task (\cref{sec:pre-task}), the two benchmarks (\cref{sec:pre-bench}), the metrics (\cref{sec:pre-metrics}), the
models and their deployment (\cref{sec:pre-models}), and the protocol and statistics used for every comparison
(\cref{sec:pre-protocol}).

\section{Task definition}
\label{sec:pre-task}
An episode gives an instruction $\ell$ that identifies one object with floor position $x^\star\in\mathbb{R}^2$. At
step $t$ the agent observes an RGB frame $I_t$, a depth image $D_t$, a pose $p_t=(x_t,\theta_t,\phi_t)$ consisting of
floor position, heading, and camera tilt, and the success flag $o_{t-1}$ of its last action. It chooses an action
$u_t\in\mathcal{U}$, where $\mathcal{U}$ contains four 0.25\,m translations (forward, backward, left, right), two
90$^\circ$ turns, and two 30$^\circ$ camera tilts. The episode is successful if
\begin{equation}
\exists\, t\le T:\ \lVert x_t-x^\star\rVert\le 1\,\text{m},\qquad T=20 .
\label{eq:success}
\end{equation}
A rejected action, for instance a translation into furniture, leaves the pose unchanged and sets $o_t=0$. The agent
never observes $x^\star$; it observes only the instruction and its own sensors.

\Cref{tab:notation} collects the symbols used in the rest of the thesis.

\begin{table}[htbp]
\centering\small
\caption{Notation.}
\label{tab:notation}
\begin{tabular*}{\linewidth}{@{\extracolsep{\fill}}ll@{}}
\toprule
Symbol & Meaning\\
\midrule
$\ell$, $s$ & instruction; target phrase extracted from it\\
$I_t$, $D_t$, $p_t$ & RGB frame, depth image, and pose $(x_t,\theta_t,\phi_t)$ at step $t$\\
$u_t$, $o_t$ & action and its success flag\\
$f$ & the frozen VLM\\
$\pi_t$, $b_t$ & presence probability and box returned by $f$ for $s$ on $I_t$\\
$a_t$ & weight of the observation of step $t$\\
$q_t$ & floor point to which a box is lifted\\
$\mu_k$, $w_k$, $V_k$ & position, evidence, and viewing poses of hypothesis $k$\\
$M_t$ & occupancy map\\
$\mathcal{B}_t$ & belief, the hypotheses together with the map\\
$\tau_\mathrm{c}$, $\tau_\mathrm{keep}$ & commit and keep thresholds\\
\bottomrule
\end{tabular*}
\end{table}

\section{Benchmarks}
\label{sec:pre-bench}
\subsection{EB-Navigation}
EB-Navigation~\cite{yang2025embodiedbench} poses the task of \cref{sec:pre-task} in single AI2-THOR
rooms~\cite{kolve2017ai2thor}. Its test set has 300 episodes in five subsets of 60, which differ only in how the
target is described or where the agent starts.
\begin{description}
\item[Base] names the target category, e.g.\ ``navigate to the AlarmClock in the room''.
\item[Common] describes the object by its use or by common sense instead of its name.
\item[Complex] embeds the target in a longer instruction with irrelevant context.
\item[Visual] describes only the appearance of the object: colour, shape, or material.
\item[Long] repeats Base with the agent turned away from the target, so the target starts behind the agent.
\end{description}
\Cref{tab:subsets} shows the five instructions of one test scene, FloorPlan13, in which the target is a toaster. The
same object is named directly, described by its use, buried in an unrelated description of the room, described only
by its colour and shape, and named again with the agent facing away from it.

\begin{table}[htbp]
\centering\small
\caption[The five EB-Navigation subsets]{The five EB-Navigation subsets, illustrated by the test instructions of
scene FloorPlan13 (target: a toaster).}
\label{tab:subsets}
\begin{tabular}{@{}lp{0.74\linewidth}@{}}
\toprule
Subset & Instruction\\
\midrule
Base & ``navigate to the Toaster in the room and be as close as possible to it''\\
Common & ``I'd like a machine that could toast some slices of bread until they are crisp. Can you navigate to that
object and stay close?''\\
Complex & ``The room is filled with the golden glow of the late afternoon sun, making it feel warm and inviting.
There's a fruit basket nearby, with bananas and apples spilling out. Please move towards the toaster. Everything feels
so peaceful right now.''\\
Visual & ``Approach the red rectangular item with black edges around the sides.''\\
Long & ``navigate to the Toaster in the room and be as close as possible to it'' (start heading rotated by
180$^\circ$)\\
\bottomrule
\end{tabular}
\end{table}

The benchmark ends an episode as soon as the agent is within 1\,m of the target, so on EB-Navigation an agent never
has to decide when to stop. All experiments use the official environment, action set, and success test. The camera
renders $500{\times}500$ RGB and depth images with a 100$^\circ$ field of view at 1.576\,m above the floor, and the
agent body is modelled as a disc of 0.2\,m radius. Translations snap to the simulator's grid, so a 0.25\,m move may
end up to a few centimetres away from its nominal position.

\subsection{RoboTHOR ObjectNav}
RoboTHOR~\cite{deitke2020robothor} poses object-goal navigation in apartments with several rooms. The validation
split used here has \nRoboN{} episodes in 15 apartments with 12 target categories. The agent uses LoCoBot actuation
with stochastic motion, 30$^\circ$ turns, and at most 500 steps. Unlike EB-Navigation, it must issue a Stop action;
an episode succeeds only if Stop is issued within 1\,m of a visible object of the target category. RoboTHOR therefore
tests two things that EB-Navigation does not: larger scenes and an explicit stop decision. Its camera renders
$640{\times}480$ images with a 63.5$^\circ$ field of view, which shows less of the room per frame than the wide
camera of EB-Navigation. \Cref{tab:bench-compare} summarizes the differences between the two benchmarks.

\begin{table}[htbp]
\centering\small
\caption{The two benchmarks.}
\label{tab:bench-compare}
\begin{tabular*}{\linewidth}{@{\extracolsep{\fill}}lll@{}}
\toprule
 & EB-Navigation & RoboTHOR ObjectNav\\
\midrule
Scenes & single iTHOR rooms & 15 multi-room apartments\\
Episodes evaluated & 300 (test) & \nRoboN{} (validation)\\
Goal description & five instruction types & 12 object categories\\
Step budget & 20 & 500\\
Translations & 0.25\,m, four directions & 0.25\,m forward, stochastic\\
Turns, tilts & 90$^\circ$, 30$^\circ$ & 30$^\circ$, 30$^\circ$\\
Camera & $500{\times}500$, 100$^\circ$ & $640{\times}480$, 63.5$^\circ$\\
End of episode & automatic on success & agent issues Stop\\
Metrics & SR & SR, SPL\\
\bottomrule
\end{tabular*}
\end{table}

\section{Evaluation metrics}
\label{sec:pre-metrics}
The main metric is the \emph{success rate} (SR), the fraction of episodes that satisfy the success condition. On
RoboTHOR, \emph{success weighted by path length} (SPL)~\cite{anderson2018evaluation} additionally rewards short paths,
\begin{equation}
\mathrm{SPL}=\frac{1}{N}\sum_{i=1}^{N} S_i\,\frac{L_i^\star}{\max(L_i,L_i^\star)},
\label{eq:spl}
\end{equation}
where $S_i$ is the success indicator of episode $i$, $L_i$ the length of the agent's path, and $L_i^\star$ the
length of the shortest path to the target. On EB-Navigation the thesis also reports the number of steps to success,
success within a step budget, and the number of VLM calls and tokens per episode.

To explain failures, each failed EB-Navigation episode is assigned the first matching outcome of a fixed taxonomy:
the target never enters the view (\emph{never seen}); the final distance is at most 1.5\,m (\emph{near miss}); the
target is out of view in the last five steps (\emph{lost after seen}); at least 30\% of the actions are rejected
(\emph{blocked}); at least three immediate reversals (\emph{oscillation}); otherwise \emph{approach failure}.
Visibility comes from the simulator's object metadata, which is used for analysis only and never given to an agent.

\section{Models and deployment}
\label{sec:pre-models}
Two open VLMs of different generations serve as the agent's model: Qwen3-VL-4B-Instruct~\cite{bai2025qwen3vl} and
Qwen2.5-VL-3B-Instruct~\cite{bai2025qwen25vl}. Both are served locally with vLLM~\cite{kwon2023vllm} at temperature 0
on one consumer GPU with 8\,GB of memory (an NVIDIA RTX~5060), which also renders the simulator. \ours uses the 4-bit
AWQ checkpoints of both models~\cite{lin2024awq}. Monocular depth for the RGB-only configuration comes from Depth
Anything V2 Metric-Indoor-Base~\cite{yang2024depthv2}, which runs on the CPU. No model is fine-tuned at any point. \Cref{tab:system} lists the hardware and software. The GPU is shared by the
VLM server and the simulator's renderer, so the model server is limited to a fraction of the 8\,GB of GPU memory.

\begin{table}[htbp]
\centering\small
\caption{Hardware and software used for all experiments.}
\label{tab:system}
\begin{tabular}{@{}lp{0.68\linewidth}@{}}
\toprule
Component & Setting\\
\midrule
GPU & one NVIDIA GeForce RTX 5060, 8\,GB, shared by the VLM server and the renderer\\
CPU and memory & 28 hardware threads, 32\,GB RAM\\
Operating system & Ubuntu 20.04\\
Simulator & AI2-THOR 5.0 with cloud rendering\\
VLM serving & vLLM 0.32, temperature 0, 4-bit AWQ weights for \ours\\
Models & Qwen3-VL-4B-Instruct, Qwen2.5-VL-3B-Instruct\\
Monocular depth & Depth Anything V2 Metric-Indoor-Base, on the CPU\\
\bottomrule
\end{tabular}
\end{table}

\section{Protocol and statistics}
\label{sec:pre-protocol}
\subsection{Development and test splits}
All design choices and parameters of \ours were selected on a separate development split and frozen before its
test runs. EmbodiedBench provides targets and instruction variants for 90 iTHOR rooms, and its test set uses 60
of them. Rerunning the benchmark's own generator on the other 30 rooms gives 150 development targets in five subsets of
30; three generator seeds give each target three start poses, 450 development episodes in all (\cref{sec:appB-dev}). Test results are reported only on the complete official splits.

For the Qwen3-VL-2B EB-prompt baseline, frequent unparseable outputs led us to stop an incomplete test run.
Before the final test comparison, we fixed a rule to select the official or a length-bounded output schema using
150 development episodes; the bounded schema was selected and the partial test run discarded. The EB baseline
uses FP8 weights, while \ours uses 4-bit AWQ weights.

\subsection{Statistical tests}
All agents run on the same 300 test episodes, so each comparison uses per-episode differences. \ours was run three
times on the EB-Navigation test set and the RGB-only configuration twice; a configuration run several times enters
with its per-episode mean success. Confidence intervals are percentile intervals of $10^4$ episode-level bootstrap
resamples~\cite{efron1993bootstrap}. $p$-values are two-sided sign-flip permutation tests with $10^6$ random sign
assignments. The seven ablation variants of \cref{sec:exp-ablation} are corrected together with the Holm
procedure~\cite{holm1979simple}, and comparisons of two single binary runs use the exact McNemar
test~\cite{mcnemar1947note}.

\subsection{Why paired tests}
All agents are evaluated on the same 300 episodes, and episodes differ greatly in difficulty: a target in plain
view at the start is solved by almost every agent, while a small target behind furniture is solved by few. An
unpaired comparison of two success rates treats this variation as noise and needs large differences to reach
significance. A paired comparison removes it: for each episode it takes the difference between the two agents'
outcomes, so only episodes on which the agents disagree carry information. The bootstrap resamples these
per-episode differences, and the sign-flip test asks how often a random assignment of signs to them produces a mean
difference as large as the observed one. When one agent was run several times, its per-episode mean success
enters the comparison, which uses all runs without counting the same episode twice. The Holm procedure controls the
probability of any false positive among the seven ablation variants at 5\%, which is stricter than testing each
variant on its own.

\section{Summary}
This chapter defined a navigation task with a 20-step budget and eight discrete actions, two benchmarks that test it
in single rooms and in apartments with an explicit stop, and a protocol in which every comparison is paired over the
same episodes and every parameter is fixed on a separate development split. The next chapter uses this setting to
diagnose why a small VLM fails when it is prompted as the agent.
